% !PRINTED 43
\weberpage{43}

%<<CONT
die Exponenten $\alpha$ in beiden dieselben sind; dann sind beide
Gruppen isomorph auf einander bezogen.

Haben also zwei Gruppen dieselben Invarianten, so können
wir diese Invarianten in beiden Gruppen zu Graden der Basiselemente
 machen, und erhalten dann diesen Fall.

Wenn umgekehrt zwei Gruppen isomorph sind, so bilden
die den Basiselementen der einen Gruppe entsprechenden Elemente
 der anderen eine Basis der letzteren, und also müssen
auch die Invarianten in beiden dieselben sein\textsuperscript{1}).

\begin{center}
§. 11.

\gesperrt{Gruppencharaktere}.
\end{center}

Es ist ein Hauptproblem in der Theorie der \gesperrt{Abel}'schen
Gruppen, alle Theiler einer \gesperrt{Abel}'schen Gruppe zu finden. Die
Lösung dieser Aufgabe wird wesentlich erleichtert durch die
Einführung des Begriffes der Gruppencharaktere, der auch sonst
in mancher Beziehung von Wichtigkeit ist.

Es sei $S$ eine \gesperrt{Abel}'sche Gruppe $n^{\text{ten}}$ Grades, deren Elemente
 durch $A$ bezeichnet werden sollen, und es seien
\weq{1}{A_1, A_2 \ldots A_\nu}
die Elemente einer Basis von $S$ von den Graden $a_1, a_2 \ldots a_\nu$.
Jedes Element $A$ ist also, und zwar nur auf eine Weise, in
der Form
\weq{2}{A = A_1^{\alpha_1} A_2^{\alpha_2} \ldots A_\nu^{\alpha_\nu}}

\textsuperscript{1}) Ueber die Theorie der \gesperrt{Abel}'schen Gruppen ist zu vergleichen:

\gesperrt{Gauss}, D\'emonstration de quelques th\'eor\`emes concernant les
périodes des classes des formes binaires du second degré. Werke,
Bd. II, S. 266.

\gesperrt{Schering}, Die Fundamentalclassen der zusammensetzbaren
arithmetischen Formen, Göttinger Abhandlungen, Bd. 14.

\gesperrt{Frobenius} und \gesperrt{Stickelberger}, Ueber Gruppen vertauschbarer
 Elemente. Crelle's Journal, Bd. 86, S. 217.

\gesperrt{Weber}, ``Theorie der \gesperrt{Abel}'schen Zahlkörper'', Acta mathematica,
Bd. 8 u. 9.

Der Begriff der Invarianten ist zuerst eingeführt in der Abhandlung
von \gesperrt{Frobenius} und \gesperrt{Stickelberger}, aber etwas anders definirt als
hier. Dieser älteren Definition, die sich bei der Anwendung auf die Kreistheilungszahlen
 minder zweckmässig erweist, ist der Verfasser in der
citirten Abhandlung in den Acta mathematica gefolgt.
%CONT>>
